%%
%% sample-sigconf.tex
%% ACM conference proceedings template – filled with paper‑specific content
%%
\documentclass[sigconf]{acmart}

\AtBeginDocument{%
  \providecommand\BibTeX{{\normalfont B\kern-0.1em{\scshape i\kern-0.1em b}\kern-0.1em\TeX}}}

\setcopyright{acmlicensed}
\copyrightyear{2026}
\acmYear{2026}
\acmDOI{XX.XXXX/XXXXXXX.XXXXXXX}

\acmConference[KDD '26]{ACM SIGKDD Conference on Knowledge Discovery and Data Mining}{August 2026}{Location TBD}
\acmISBN{978‑1‑4503‑XXXX‑X/26/08}

\citestyle{acmauthoryear}

% ------------------------------------------------------------------
\begin{document}

\title[Context is Everything]{Context is Everything: Boosting Cricket Dismissal Prediction with Large Language Models}

\author{Mustafa Ashfaq}
\orcid{0000-000X-XXXX-XXXX}
\affiliation{%
  \institution{LUMS}
  \city{Lahore}
  \country{Pakistan}}
\email{25100216@lums.edu.pk}

\author{Essa Shahid Arshad}
\orcid{0000-000X-XXXX-XXXX}
\affiliation{%
  \institution{LUMS}
  \city{Lahore}
  \country{Pakistan}}
\email{25100297@lums.edu.pk}

\author{Zafar Ayyub Qazi}
\orcid{0000-000X-XXXX-XXXX}
\affiliation{%
  \institution{LUMS}
  \city{Lahore}
  \country{Pakistan}}
\email{zafar.qazi@lums.edu.pk}

\author{Ihsan Ayyub Qazi}
\orcid{0000-000X-XXXX-XXXX}
\affiliation{%
  \institution{LUMS}
  \city{Lahore}
  \country{Pakistan}}
\email{ihsan.qazi@lums.edu.pk}

% ------------------------------------------------------------------
\begin{abstract}
Cricket ranks as the world's second-most popular sport with approximately 2.5 billion fans, yet predicting how a batsman will be dismissed remains exceptionally challenging. Each dismissal type (bowled, caught, LBW, stumped, run‑out, etc.) is shaped by layered contextual factors such as venue, pitch, weather, opponent, toss, bowler style, field setup, and match state. Traditional predictive models, including Bayesian survival analysis, hidden Markov models, and random forests, often overlook critical context due to data sparsity, high-dimensionality, and operational constraints, leading to limited scalability and frequent retraining needs. While Large Language Models (LLMs) offer the ability to tailor predictions to each player through in-context learning, it remains unknown which contextual inputs maximize prediction accuracy. In this paper, we systematically evaluate LLMs for batsman dismissal prediction by evaluating the impact of different contextual inputs. Through comprehensive ablation experiments using structured (career statistics, venue and opponent history) and unstructured (recent match commentary) contextual inputs across multiple LLM variants, we find that including venue and opponent context elevates top-1 accuracy from 43\% (career stats alone) to 57.5\%, and top-2 accuracy from 57\% to 79.2\%. Among tested models, GPT-4.1-mini offers the optimal trade-off between accuracy and efficiency for practical deployment. Our results establish which contextual factors most enhance predictive performance, provide a scalable and interpretable LLM-based framework for cricket analytics, and outline clear directions for future work in sports prediction.
\end{abstract}

% ------------------------------------------------------------------
\begin{CCSXML}
<ccs2012>
   <concept>
      <concept_id>10010147.10010257.10010258.10010261</concept_id>
      <concept_desc>Computing methodologies~Machine learning</concept_desc>
      <concept_significance>500</concept_significance>
   </concept>
   <concept>
      <concept_id>10010405.10010444.10010449</concept_id>
      <concept_desc>Applied computing~Sports</concept_desc>
      <concept_significance>300</concept_significance>
   </concept>
</ccs2012>
\end{CCSXML}

\ccsdesc[500]{Computing methodologies~Machine learning}
\ccsdesc[300]{Applied computing~Sports}

\keywords{cricket analytics, large language models, contextual modeling, dismissal prediction, sports AI, ablation study, structured prediction}

% ------------------------------------------------------------------
\maketitle

\section{Introduction}

\begin{itemize}
\item Cricket's global reach and popularity (2.5 billion fans worldwide)
\item Challenge: Predicting batsman dismissals is exceptionally difficult due to complex contextual factors
\item Current state: Traditional ML approaches (Random Forest, Naive Bayes, HMMs) have fundamental limitations
\item Empirical evidence: Traditional models achieve only 22-31\% accuracy vs our LLM approach (57.5\%)
\item Problem complexity: 8 dismissal types shaped by venue, opponent, pitch, weather, bowling style, field setup
\item Research gap: Unknown which contextual inputs maximize LLM prediction accuracy
\item Our contribution: First systematic evaluation of LLMs for cricket dismissal prediction with comprehensive ablation study
\item Key finding: Venue + opponent context improves accuracy from 43\% to 57.5\% (top-1) and 57\% to 79.2\% (top-2)
\item Practical impact: GPT-4.1-mini offers optimal accuracy-cost tradeoff for real-world deployment
\item Broader significance: Establishes scalable framework for sports analytics and structured prediction tasks
\end{itemize}

\section{Related Work}

\subsection{Sports Analytics and Machine Learning}

\begin{itemize}
\item Traditional cricket analytics: Statistical models, regression analysis, survival analysis approaches
\item Existing dismissal prediction: Rule-based systems, decision trees, ensemble methods
\item Performance limitations: Poor generalization, high retraining frequency, context blindness
\item Sports prediction challenges: Data sparsity, temporal dependencies, situational complexity
\end{itemize}

\subsection{Large Language Models in Sports}

\begin{itemize}
\item Emerging applications: Sports commentary generation, strategy analysis, outcome prediction
\item Advantages: Natural language processing, contextual understanding, in-context learning
\item Sports-specific challenges: Structured data integration, domain adaptation, evaluation metrics
\item Gap: Systematic evaluation of contextual inputs for sports prediction tasks
\end{itemize}

\subsection{Contextual Learning and Structured Prediction}

\begin{itemize}
\item In-context learning capabilities: Few-shot learning, prompt engineering, chain-of-thought reasoning
\item Structured prediction: Multi-class classification, sequential decision making, constraint satisfaction
\item Context incorporation: Textual context, temporal patterns, domain-specific knowledge
\item Evaluation challenges: Metrics design, baseline establishment, statistical significance testing
\end{itemize}

\section{Problem Formulation}

\begin{itemize}
\item Task definition: 8-class cricket dismissal prediction (bowled, caught-fielder, caught-wicketkeeper, stumped, LBW, hit-wicket, caught-and-bowled, not-out)
\item Input space: Structured data (career stats, venue history, opponent records) + Unstructured data (match commentary, conditions)
\item Challenge characterization: Data sparsity (empirically shown: sparsity ratio = 1.0), high-dimensionality, temporal dependencies
\item Evaluation framework: Top-1 and Top-2 accuracy metrics, statistical significance testing (McNemar's test)
\item Traditional model limitations: Cannot process textual context, require manual feature engineering, poor scalability
\item Research questions: (1) Which contextual inputs maximize prediction accuracy? (2) How do different LLM variants compare? (3) What is the optimal accuracy-cost tradeoff?
\end{itemize}

\section{Methodology}

\subsection{Data Collection Pipeline}

\begin{itemize}
\item ESPNCricinfo scraping: Automated data extraction from 3,000+ ODI matches
\item Temporal coverage: Multi-year span capturing diverse playing conditions and rule changes
\item Data types: Match scorecards, ball-by-ball commentary, player statistics, venue information
\item Player selection: 7 international players across multiple teams and playing styles
\item Data volume: 1,700+ processed commentary files, 600+ aggregated player datasets, 100+ cricket grounds
\end{itemize}

\subsection{Contextual Input Framework}

\begin{itemize}
\item Structured contexts:
\begin{itemize}
\item Career statistics: Batting averages, strike rates, dismissal patterns by bowler type
\item Venue-specific performance: Ground-wise statistics, pitch condition adaptation
\item Opponent-specific analysis: Team-wise historical performance, matchup patterns
\item Temporal patterns: Recent form, seasonal variations, career trajectory
\end{itemize}
\item Unstructured contexts:
\begin{itemize}
\item Ball-by-ball commentary: Real-time match descriptions, player behavior analysis
\item Match conditions: Pitch reports, weather conditions, venue characteristics
\item Strategic context: Field placements, bowling changes, pressure situations
\end{itemize}
\end{itemize}

\subsection{Ablation Study Design}

\begin{itemize}
\item Systematic evaluation: 9 contextual input combinations tested
\item Control variables: Player identity, LLM model, evaluation protocol maintained constant
\item Progressive addition: Baseline (career only) → +Venue → +Opponent → +Commentary
\item Experimental matrix: All combinations of structured and unstructured inputs evaluated
\item Reproducibility: Automated experimental framework with consistent prompt engineering
\end{itemize}

\section{Experimental Setup}

\subsection{Model Selection and Configuration}

\begin{itemize}
\item LLM variants tested: GPT-4.1, GPT-4.1-mini, GPT-4.1-nano, o4-mini
\item Prompt engineering: Chain-of-thought reasoning, systematic analysis steps, consistent output formatting
\item API configuration: Temperature settings, token limits, retry mechanisms for reliability
\item Cost analysis: Computational overhead, API costs, deployment considerations
\end{itemize}

\subsection{Baseline Comparisons}

\begin{itemize}
\item Traditional ML models: Random Forest (24.3\% accuracy), Naive Bayes (22.3\%), Logistic Regression (26.2\%), SVM (31.1\%)
\item Feature engineering: Manual conversion of cricket statistics to numerical features
\item Cross-validation: 5-fold CV with high variance (±4-6\% standard deviation)
\item Limitations demonstrated: Data sparsity issues, inability to process textual context, poor generalization
\end{itemize}

\subsection{Evaluation Protocol}

\begin{itemize}
\item Test dataset: 501 dismissal instances across 7 players
\item Metrics: Top-1 and Top-2 prediction accuracy
\item Statistical validation: McNemar's test for pairwise significance testing
\item Cross-validation: Player-wise and temporal splits to ensure robustness
\item Confidence intervals: Bootstrap sampling for uncertainty quantification
\end{itemize}

\section{Results}

\subsection{Main Findings}

\begin{itemize}
\item Performance improvement: 43\% → 57.5\% (top-1 accuracy) with venue + opponent context
\item Top-2 accuracy enhancement: 57\% → 79.2\% with optimal contextual inputs
\item Best performing model: o4-mini achieves highest accuracy (57.5\% top-1, 79.2\% top-2)
\item Traditional model comparison: LLMs outperform traditional ML by 26.4\% absolute improvement
\item Statistical significance: McNemar's test confirms significant differences (p < 0.01) between context combinations
\end{itemize}

\subsection{Ablation Analysis}

\begin{itemize}
\item Progressive context addition results:
\begin{itemize}
\item Baseline (career stats only): 43\% top-1 accuracy
\item +Venue context: 51.3\% (+8.3\% improvement)
\item +Opponent context: 57.5\% (+6.2\% additional improvement)
\item +Commentary: Marginal gains (varies by model)
\end{itemize}
\item Context effectiveness ranking: Venue + Opponent > Venue only > Opponent only > Commentary only
\item Model consistency: Results stable across different LLM variants with similar patterns
\end{itemize}

\subsection{Model Comparison Analysis}

\begin{itemize}
\item Performance ranking: o4-mini (57.5\%) > GPT-4.1-mini (56.3\%) > GPT-4.1 (55.3\%) > GPT-4.1-nano (51.3\%)
\item Cost-accuracy tradeoff: GPT-4.1-mini offers optimal balance for practical deployment
\item Consistency analysis: Lower variance in LLM predictions compared to traditional models
\item Computational efficiency: Model-specific analysis of inference time and resource requirements
\end{itemize}

\subsection{Statistical Significance Testing}

\begin{itemize}
\item McNemar's test results: Pairwise comparisons between all contextual configurations
\item Significant improvements: Venue + opponent vs baseline (p < 0.001)
\item Non-significant differences: Commentary addition shows inconsistent benefits
\item Confidence intervals: Bootstrap analysis confirms robustness of main findings
\end{itemize}

\section{Discussion}

\subsection{Why Contextual Inputs Matter}

\begin{itemize}
\item Venue effects: Ground dimensions, pitch characteristics, weather patterns significantly impact dismissal patterns
\item Opponent strategies: Bowling styles, field placements, tactical approaches vary by team
\item Commentary limitations: Real-time context vs historical data tradeoffs, inconsistent information quality
\item Empirical validation: Venue + opponent context captures crucial predictive signals missed by career statistics alone
\end{itemize}

\subsection{Traditional Model Limitations}

\begin{itemize}
\item Data sparsity: Maximum sparsity ratio (1.0) prevents effective pattern learning
\item Context blindness: Cannot process crucial textual information (match commentary, venue descriptions, strategic insights)
\item Feature engineering burden: Manual conversion loses important contextual nuances
\item Scalability issues: Require complete retraining for new venues, opponents, or rule changes
\item Operational constraints: Multiple specialized libraries, extensive hyperparameter tuning, domain expertise requirements
\end{itemize}

\subsection{LLM Advantages}

\begin{itemize}
\item Natural context processing: Direct integration of textual match information without feature engineering
\item In-context learning: Adaptation to new scenarios without retraining
\item Chain-of-thought reasoning: Systematic analysis incorporating multiple information sources
\item Scalability: Single framework handles diverse input types and contexts
\item Practical deployment: Minimal infrastructure requirements compared to traditional ML pipelines
\end{itemize}

\subsection{Implications for Sports Analytics}

\begin{itemize}
\item Broader applications: Framework applicable to other sports prediction tasks
\item Real-world deployment: Broadcasting analysis, betting markets, coaching insights
\item Methodological contributions: Best practices for LLM evaluation in sports analytics
\item Cost considerations: Balance between prediction accuracy and computational overhead
\end{itemize}

\section{Limitations and Future Work}

\begin{itemize}
\item Dataset scope: Focus on ODI cricket, potential for extension to T20 and Test formats
\item Player selection: Limited to 7 players, expansion needed for broader generalization
\item Temporal factors: Recent vs historical data weighting requires further investigation
\item Commentary quality: Dependency on availability and quality of textual match descriptions
\item Real-time prediction: Extension to live match scenarios with streaming data integration
\item Multi-modal inputs: Integration of video analysis, player tracking, and other data sources
\end{itemize}

\section{Conclusion}

\begin{itemize}
\item Research contributions: First systematic evaluation of LLMs for cricket dismissal prediction
\item Key findings: Venue + opponent context provides 26.4\% improvement over traditional models
\item Methodological insights: Structured contextual input evaluation framework for sports analytics
\item Practical impact: GPT-4.1-mini offers optimal accuracy-cost tradeoff for deployment
\item Broader significance: Advances application of LLMs to structured prediction tasks in sports
\item Future directions: Extension to other sports, real-time prediction, multi-modal integration
\end{itemize}

\begin{acks}
We thank the cricket analytics community for valuable feedback and ESPNCricinfo for data availability. This research was supported by LUMS research facilities and computational resources.
\end{acks}

\bibliographystyle{ACM-Reference-Format}
\bibliography{cricket-references}

\end{document}
